\documentclass[11pt]{article}

% --------------------------------------------------
% Packages
% --------------------------------------------------
\usepackage[T1]{fontenc}
\usepackage{lmodern}
\usepackage{geometry}
\usepackage{microtype}
\usepackage{setspace}
\usepackage{amsmath,amssymb,amsthm,mathtools}
\usepackage{csquotes}
\usepackage{hyperref}
\usepackage{booktabs}

\geometry{margin=1in}
\setstretch{1.15}

% --------------------------------------------------
% Theorem environments
% --------------------------------------------------
\newtheorem{definition}{Definition}
\newtheorem{proposition}{Proposition}
\newtheorem{remark}{Remark}

% --------------------------------------------------
% Macros
% --------------------------------------------------
\newcommand{\doop}{\mathrm{do}}
\newcommand{\Hist}{\mathcal{H}}
\newcommand{\State}{\mathcal{S}}
\newcommand{\Alg}{\mathcal{A}}

% --------------------------------------------------
% Title
% --------------------------------------------------
\title{From Inspection to Intervention:\\
Causal Mechanisms and Algorithmic Understanding in Neural Networks}

\author{Flyxion}

\date{December 2025}

\begin{document}
\maketitle

% ==================================================
\begin{abstract}
% ==================================================

Recent advances in mechanistic interpretability have shifted attention from
descriptive inspection of neural networks toward causal analysis of their
internal operations. This essay synthesizes a causal paradigm in which
mechanistic understanding is defined not by access to internal states, but by
the ability to perform interventions, evaluate counterfactuals, and recover
algorithmic structure through interventional equivalence. We examine the
progression from activation-level analysis to causal mediation and abstraction,
and we argue that genuine understanding requires sensitivity to history,
irreversibility, and responsibility rather than mere correlational control. The
resulting framework reframes interpretability as a problem of reverse-engineering
algorithms implemented by dynamical systems under constraint.

\end{abstract}

% ==================================================
\section{Introduction}
% ==================================================

The problem of understanding neural networks has long been framed as one of
visibility. If the internal weights, activations, or attention patterns could be
made sufficiently transparent, then the system’s behavior would become
intelligible. This intuition has motivated a wide range of visualization and
probing techniques that aim to ``open the black box.''

However, visibility alone does not constitute understanding. Observing that a
component activates in the presence of a behavior does not establish that the
component is responsible for that behavior, nor does it explain how the behavior
would change under alternative conditions. The gap between observation and
explanation motivates a more demanding notion of mechanistic understanding, one
rooted in causality rather than correlation.

This essay develops such a notion. Mechanistic understanding is taken to mean
the ability to identify internal components whose manipulation predictably and
systematically alters behavior in a way that corresponds to an intelligible
algorithm. On this view, interpretability becomes a causal science.

% ==================================================
\section{What It Means to Be Mechanistic}
% ==================================================

The term ``mechanistic'' is often used ambiguously. In its strongest technical
sense, a mechanistic explanation specifies not merely which parts are present in
a system, but how those parts causally interact to produce behavior.

\begin{definition}[Mechanistic Understanding]
A system is mechanistically understood with respect to a behavior if one can
identify a set of internal variables such that interventions on those variables
produce predictable, structured changes in the behavior, consistent with an
underlying algorithmic description.
\end{definition}

This definition excludes purely observational accounts. A neuron that reliably
activates during correct answers may be diagnostically useful, but it is not
mechanistically explanatory unless altering its value changes the outcome in a
lawful way.

% ==================================================
\section{Causality and Intervention}
% ==================================================

Causality enters the picture through intervention. An intervention sets an
internal variable to a value it would not ordinarily take, breaking its normal
dependencies and revealing its downstream effects.

Formally, an intervention replaces the natural evolution of a variable $X$ with
an assignment $\doop(X = x)$, holding fixed the rest of the system’s causal
structure. The resulting behavior distinguishes genuine causal influence from
mere correlation.

Counterfactual reasoning extends this idea. Rather than asking what happens when
a variable is zeroed or scaled, one asks what would have happened had the system
internally evolved as if it were processing a different input. Such
counterfactuals probe whether a component carries information that is
functionally responsible for a computation.

% ==================================================
\section{Activation-Level Control}
% ==================================================

One of the earliest demonstrations of causal influence comes from activation
steering. By adding a vector to internal activations, researchers can bias a
model toward or away from certain behaviors. When steering vectors are derived
from differences in mean activation between contrasting behaviors, they reveal
that behavioral distinctions often correspond to approximately linear
directions in representation space.

Despite its usefulness, activation steering remains limited. It demonstrates
that internal representations are manipulable, but it does not by itself specify
what computation is being performed. Control without explanation is not yet
mechanistic understanding.

% ==================================================
\section{Causal Mediation and Information Flow}
% ==================================================

Causal mediation analysis refines intervention by decomposing the total effect of
an input on an output into contributions carried by intermediate variables.
Rather than asking whether a component matters, mediation asks how and where
information flows.

In causal tracing experiments, noise is injected to degrade a model’s
performance, after which selected components are restored to their clean values.
Components whose restoration recovers correct behavior are identified as
mediators. This procedure reveals structured pathways by which information is
retrieved, transformed, and routed through the network.

Such analyses produce causal stories rather than static maps. They show not only
where information is represented, but how it moves.

% ==================================================
\section{Causal Abstraction and Algorithmic Structure}
% ==================================================

The strongest form of mechanistic interpretability arises when neural activity
can be related to a higher-level algorithm through causal abstraction. Both the
neural network and the algorithm are treated as causal models, and the goal is
to determine whether the algorithm faithfully abstracts the network’s behavior
under intervention.

\begin{definition}[Causal Abstraction]
A high-level model $\Alg$ is a causal abstraction of a low-level system if there
exists a mapping between their variables such that corresponding interventions
produce equivalent effects on outputs.
\end{definition}

Interchange interventions play a central role here. Instead of modifying a
variable arbitrarily, one replaces its value with the value it would have taken
under a different input. If such substitutions produce output changes that match
the predictions of the abstract algorithm, the abstraction is validated.

This criterion is stringent. It requires not just similarity of behavior, but
interventional equivalence.

% ==================================================
\section{Events, History, and Irreversibility}
% ==================================================

Causal abstraction implicitly raises questions about the role of history. Two
internal states may be numerically identical yet differ in how they were
produced. If their downstream effects diverge under intervention, then the
mechanism is history-sensitive.

This suggests that mechanistic explanation may require attention not only to
states, but to events: transformations that irreversibly constrain future
behavior. Interventions that cannot be undone without loss of explanatory power
signal the presence of genuine historical commitment within the system.

% ==================================================
\section{Responsibility and Constraint}
% ==================================================

Not all mechanisms act by producing outcomes. Some function by suppressing,
filtering, or excluding possibilities. Such negative or constraining roles are
often invisible to activation-based analyses, yet they may be causally decisive.

A component may bear responsibility for a behavior not by optimizing for it, but
by preventing alternatives. Recognizing such mechanisms requires a notion of
causality that treats absence and inhibition as first-class explanatory factors.

% ==================================================
\section{Limits of Understanding}
% ==================================================

Successful interventions do not automatically constitute understanding. A
collection of techniques may enable control without yielding a coherent
explanatory model. Conversely, an elegant abstraction may fail outside a narrow
regime of validity.

Mechanistic understanding therefore admits of degrees. It improves as
abstractions become more stable under distribution shift, more sensitive to
history, and more faithful under counterfactual substitution.

% ==================================================
\section{Conclusion}
% ==================================================

Mechanistic interpretability, when grounded in causality, becomes the study of
which internal structures are responsible for behavior and which algorithms they
realize. This shift moves the field beyond visualization and toward
reverse-engineering.

The central lesson is that understanding is not achieved by looking, but by
intervening; not by correlating, but by substituting; and not by controlling
outputs alone, but by reconstructing the causal logic that produces them. Only
then can claims about what a neural network is doing rise from metaphor to
mechanism.

\end{document}
